\section{Detectors exclude linear \texorpdfstring{\(c_0\)}{c0}}\label{sec:no-c0}

A linear copy of \(c_0\) would provide vectors bounded away from zero
whose finite signed sums are uniformly bounded. We first show that
one class of tests detects a subsequence of these vectors. We then
identify the additional property needed for a contradiction: the tests'
finite signed sums must have one Lipschitz bound. The rest of the
section obtains this property by modifying the detectors and separating
or nesting their supports. Supports always refer to the concrete
generated functions on \(B\).

\subsection{From bounded signed sums to individual detectors}

The extraction has two stages. The outer square sum prevents detection
from escaping through different seminorms; within one seminorm, the
exponents approaching one prevent it from spreading over arbitrarily
many small scalar pairings.

\begin{lemma}[One class detects a subsequence]\label{lem:detectors}
Suppose \((m_i)\subset Z\) satisfies, for some \(r,M>0\),
\begin{equation}
 \|m_i\|_Z\geq r,\qquad
 \left\|\sum_{i\in H}\epsilon_i m_i\right\|_Z\leq M
 \quad\text{for all finite \(H\) and all signs \(\epsilon_i\)}. \label{eq:bounded-signed}
\end{equation}
Then \(u(m_i)\to0\) for every \(u\in Z^*\). After a subsequence,
there are tests \(F_i\) in one listed class \(\mathcal L\) and a
constant \(\delta>0\) with
\begin{equation}
                         |\langle F_i,m_i\rangle|\geq\delta.
                                                               \label{eq:detected}
\end{equation}
\end{lemma}

\begin{proof}
For any fixed \(u\in Z^*\), choosing signs matching \(u(m_i)\) in
\eqref{eq:bounded-signed} gives
\[
                    \sum_{i\in H}|u(m_i)|\leq M\|u\|
\]
for every finite \(H\). Thus \(\sum_i|u(m_i)|<\infty\), and in
particular \(u(m_i)\to0\). Also \(\|m_i\|\leq M\).

We first prove that one fixed seminorm stays away from zero on a
subsequence. Otherwise \(\sigma_d(m_i)\to0\) for every \(d\).
The vectors
\[
               v_i=(\sqrt{\alpha_d}\sigma_d(m_i))_{d\geq1}
\]
then converge coordinatewise to zero in \(\ell_2\), while
\(r\leq\|v_i\|_2=\|m_i\|_Z\leq M\). A direct gliding-hump selection
gives disjoint finite intervals \(A_i\) of indices, after a
subsequence, such that
\begin{equation}
                 \sum_{d\in A_i}\alpha_d\sigma_d(m_i)^2
                            \geq r^2/4.                      \label{eq:hump-mass}
\end{equation}
Indeed, once the endpoint of the preceding interval is fixed,
coordinatewise convergence makes the squared mass before that
endpoint less than \(r^2/4\) for a sufficiently late \(v_i\).
Then choose a new finite endpoint so that the squared tail mass is
less than \(r^2/4\). The mass in between is at least \(r^2/2\),
which implies \eqref{eq:hump-mass}.

The function \(\sigma_d^2\) is even and convex, because a seminorm
is convex and squaring is convex and increasing on \([0,\infty)\).
For fixed \(y,m\),
Jensen's inequality applied to \(y+m\) and \(-y+m\) gives
\begin{equation}
       \frac{\sigma_d(y+m)^2+\sigma_d(y-m)^2}{2}
                        \geq\sigma_d(m)^2.                  \label{eq:seminorm-convexity}
\end{equation}
For independent random signs and a finite \(H\), average first over
the sign of \(m_i\) and use \eqref{eq:seminorm-convexity}. Thus
\[
 \mathbb E\,\sigma_d\!\left(\sum_{k\in H}\epsilon_km_k\right)^2
                         \geq\sigma_d(m_i)^2\quad(i\in H).
\]
The exact norm formula, \eqref{eq:bounded-signed}, the disjointness of the \(A_i\), and
\eqref{eq:hump-mass} imply
\[
 M^2\geq\mathbb E\left\|\sum_{k\in H}\epsilon_km_k\right\|_Z^2
 \geq\sum_{i\in H}\sum_{d\in A_i}\alpha_d\sigma_d(m_i)^2
 \geq |H|r^2/4,
\]
a contradiction for large \(H\). Hence some
\(d=(\mathcal L,C,p)\) has \(\sigma_d(m_i)\geq a>0\) on a subsequence.

Choose a listed exponent \(p'=1+1/k'<p\), with the same
\(\mathcal L,C\), and call its triple \(d'\). Its seminorm is bounded
on these vectors by
\[
                  \sigma_{d'}(m_i)\leq
                  \alpha_{d'}^{-1/2}M=:B_0.
\]
For each \(i\), choose a finite admissible list for \(d\) whose
\(\ell_p\) value on \(m_i\) is at least \(a/2\). Put
\(x_\nu=|\langle u_\nu,m_i\rangle|\). The same list is admissible
for \(d'\), because only the exponent changed. Consequently
\[
 (a/2)^p\leq\sum_\nu x_\nu^p
 \leq(\max_\nu x_\nu)^{p-p'}\sum_\nu x_\nu^{p'}
 \leq(\max_\nu x_\nu)^{p-p'}B_0^{p'}.
\]
The list is nonempty and \(B_0>0\). One of its tests therefore has
evaluation at least
\[
             \delta=\left(\frac{(a/2)^p}{B_0^{p'}}\right)^{1/(p-p')}>0.
\]
Choosing that test as \(F_i\) proves \eqref{eq:detected}.
\end{proof}

\subsection{The growth forced by compatible detectors}

Call a sequence of tests \emph{compatible} if all its finite signed
sums have a common Lipschitz bound. The norm was defined using precisely
such lists. Once compatible tests detect the vectors, a single
seminorm forces the following growth.

\begin{lemma}[Growth from compatible detectors]\label{lem:detector-growth}
Let \((m_i)\subset Z\), and let \(F_i\) belong to one listed class
\(\mathcal L\). Suppose there are \(\delta>0\) and an integer
\(C\geq1\) such that
\[
 |\langle F_i,m_i\rangle|\geq\delta,
 \qquad
 \Lip\!\left(\sum_{i\in H}\epsilon_iF_i\right)\leq C
\]
for every \(i\), every finite \(H\), and all signs \(\epsilon_i\).
For the listed triple \(d=(\mathcal L,C,2)\), every nonempty finite
\(H\) satisfies
\begin{equation}
 \max_{\epsilon_i\in\{-1,1\}}
 \left\|\sum_{i\in H}\epsilon_i m_i\right\|_Z
 \geq\sqrt{\alpha_d}\,\delta\sqrt{|H|}.
 \label{eq:detector-growth}
\end{equation}
\end{lemma}

\begin{proof}
The list \((F_k)_{k\in H}\) is admissible for \(d\), including when
labels repeat. For independent real random signs, the exact list
formula in Lemma~\ref{lem:duality} gives
\begin{align*}
 \mathbb E\,\sigma_d\!\left(\sum_{i\in H}\epsilon_i m_i\right)^2
 &\geq
 \mathbb E\sum_{k\in H}
       \left|\sum_{i\in H}\epsilon_i
                    \langle F_k,m_i\rangle\right|^2\\
 &=\sum_{k\in H}\sum_{i\in H}|\langle F_k,m_i\rangle|^2
 \geq |H|\delta^2.
\end{align*}
The equality uses cancellation of cross terms under the independent
signs; no cross-pairing is assumed to vanish. By
\eqref{eq:completed-norm},
\[
 \max_{\epsilon_i\in\{-1,1\}}
       \left\|\sum_{i\in H}\epsilon_i m_i\right\|_Z^2
 \geq\mathbb E\left\|\sum_{i\in H}\epsilon_i m_i\right\|_Z^2
 \geq\alpha_d|H|\delta^2.
\]
Taking square roots proves \eqref{eq:detector-growth}.
\end{proof}

It remains to make the detectors from Lemma~\ref{lem:detectors}
compatible without losing their nonzero pairings. We will use two
support configurations: disjoint supports, or later supports contained
in cylinders where earlier functions are constant.

\subsection{Support conditions for compatibility}

Localizing detecting Lipschitz functions and combining them with a
common Lipschitz bound is also central to Kalton's arguments for
metric free spaces \cite[Lemma~4.5 and proof of Theorem~4.6]{Kalton04}.
Here the supports are finite-coordinate cylinders and tail regions,
and their compatibility must be retained within the selected test
classes defining \(Z\).

\begin{lemma}[Compatible supports]\label{lem:compatible}
Let \(B\) be a convex subset of a normed space. Let \(u_i,b_i\)
be continuous real functions on \(B\), with \(b_i\) Lipschitz, and
put \(F_i=u_i b_i\). Suppose
\(\Lip(F_i)\leq L\) for every \(i\).
\begin{enumerate}
\item If the closed relative supports \(\supp_B F_i\) are pairwise
disjoint, every finite signed sum of the \(F_i\) has Lipschitz
constant at most \(L\).
\item Suppose that for each \(i\) there is a relatively open set
\(C_i\subset B\) and a scalar \(e_i\) such that \(u_i=e_i\) on
\(C_i\) and \(\supp_B u_k\subset C_i\) for every \(k>i\).
If \(\sum_i|e_i|\Lip(b_i)<\infty\), every finite signed sum of the
\(F_i\) has Lipschitz constant at most
\[
                         L+\sum_i|e_i|\Lip(b_i).
\]
If \(b_i=1\) for every \(i\), the summability condition is automatic.
\end{enumerate}
\end{lemma}

\begin{proof}
Fix a finite set of indices and signs. In the first case, at a given
point at most one of the corresponding closed supports contains that
point. All other summands vanish on a relative neighborhood, so the
sum is locally either zero or a single signed \(F_i\), with bound
\(L\).

In the second case, choose at a given point the largest selected
index \(j\) whose \(u_j\)-support contains the point. If there is no
such index, every summand vanishes on a common neighborhood, because
there are only finitely many closed supports. Otherwise all selected
later functions vanish on a neighborhood. The point lies in \(C_i\)
for each selected \(i<j\), so on a smaller neighborhood all those
\(u_i\) equal \(e_i\). On that neighborhood the sum is
\[
                \epsilon_jF_j+\sum_{\substack{i<j\\i\text{ selected}}}
                                  \epsilon_i e_i b_i,
\]
and has the asserted local Lipschitz bound.

To pass from the uniform local bound to a global bound, take any two
points of \(B\). The segment between them is contained in \(B\).
Its compact parameter interval has a finite cover by the relative
neighborhoods just obtained and hence a Lebesgue number. Subdivide
the segment finely enough that each subsegment lies in one such
neighborhood, apply the local estimate to each subsegment, and add.
The lengths of the collinear subsegments sum to the distance of the
endpoints. This proves both global assertions.
\end{proof}

\begin{lemma}[Tail separation]\label{lem:tail-separation}
Suppose \(F_i\) belong to one listed class and have representations
with bounded band indices. Suppose \(m_i\in Z\) and
\begin{equation}
 |\langle F_i,m_i\rangle|\geq\delta>0,\qquad
\supp_B F_i\subset\{\tau_{M_i}\geq c_*\},
 \quad M_i\in\N,\quad M_i\to\infty,\quad c_*>0.             \label{eq:tail-hypotheses}
\end{equation}
Then a subsequence has modified tests in one class \(\mathcal V_H\),
with detection at least \(\delta/2\), whose relative supports are
pairwise disjoint.
\end{lemma}

\begin{proof}
Choose a dyadic \(c=2^{-a}\) with \(2c<c_*\). If the original class
is \(\mathcal T_h\), view a test as a cylindrical factor with no tail
factors; if it is \(\mathcal V_h\), keep its chosen product
representation. A single integer \(H\), larger than \(h\), the
bounded band indices, and \(a\), can be chosen so that both \(F_i\)
and \(F_i b_{c,N}\) belong to \(\mathcal V_H\) for all \(i,N\).
There is room for one extra tail factor.

For each fixed \(i\), \(b_{c,N}(s)\to1\) at every \(s\in B\), since
\(\tau_N(s)\to0\). Lemma~\ref{lem:duality}, used in this one class,
therefore gives
\(\langle F_i b_{c,N},m_i\rangle\to\langle F_i,m_i\rangle\).
Choose \(N\) large enough that the difference has modulus less than
\(\delta/2\); every sufficiently large \(N\) then works. This choice
is made for the fixed vector \(m_i\).

Select indices recursively. After choosing an index \(i\) and such
an \(N_i\), take the next original index so far out that its \(M_k\)
exceeds all previously selected \(N_i\), which is possible because
\(M_k\to\infty\). Then choose its own \(N_k\) as above. For selected
\(k>i\), the new \(i\)-th support requires \(\tau_{N_i}\leq2c\)
by \eqref{eq:tail-support}. The new \(k\)-th support requires
\(\tau_{M_k}\geq c_*\) by \eqref{eq:tail-hypotheses}, hence
\(\tau_{N_i}\geq\tau_{M_k}\geq c_*>2c\). The supports are disjoint,
and the detection bound follows from the choice of \(N_i\).
\end{proof}

\subsection{Selecting compatible detectors}

The next proposition isolates the geometric part of the argument. It
uses only weak nullity of the vectors, one fixed test class, and a
uniform lower bound on the pairings. The bounded signed sums will enter
again only when we apply Lemma~\ref{lem:detector-growth}.

\begin{proposition}[Selection of compatible detectors]
\label{prop:compatible-detectors}
Let \((m_i)\subset Z\) be weakly null, meaning
\(u(m_i)\to0\) for every \(u\in Z^*\). Suppose \(F_i\) belong to one
listed class \(\mathcal L\) and
\begin{equation}
                  |\langle F_i,m_i\rangle|\geq\delta>0
                  \qquad(i\geq1).
                  \label{eq:selection-detection}
\end{equation}
There are a subsequence \((m_{i_k})\), a listed class
\(\mathcal L_1\), tests \(\widetilde F_k\in\mathcal L_1\),
\(\delta_1>0\), and an integer \(C_1\geq1\) such that
\[
 |\langle\widetilde F_k,m_{i_k}\rangle|\geq\delta_1,
 \qquad
 \Lip\!\left(\sum_{k\in H}\epsilon_k\widetilde F_k\right)\leq C_1
\]
for every \(k\), every finite \(H\), and all signs \(\epsilon_k\).
\end{proposition}

\begin{proof}
We pass to subsequences and relabel their indices throughout the proof.
Every fixed test pairs to zero in the limit, by weak nullity and
Lemma~\ref{lem:duality}.

Fix a concrete representation of each detector \(F_i\). If its class
is \(\mathcal T_h\), take \(F_i=u_i\) and no tail factor. If its class
is \(\mathcal V_h\), write \(F_i=u_i b_i\), where \(u_i\) is its
chosen cylindrical factor and \(b_i\) is its product of tail factors.
Let \(\gamma_i=(\lambda_i,j_i,\xi_i)\) be the chosen label.
After subsequences, either the labels are pairwise distinct or one
label is repeated throughout. In the distinct case the bands have a
bounded subsequence or a subsequence tending to infinity. We handle
these three possibilities separately.

\paragraph{Distinct labels with bounded bands.}
Every block \(\Gamma_\lambda\) is finite, so pairwise distinct labels
force \(\lambda_i\) to leave every finite set. Since
\(\{\lambda:J_\lambda\leq J\}\) is finite for every \(J\), we have
\(J_{\lambda_i}\to\infty\). Pass to a further subsequence with
one fixed band \(j\). Eventually \(j<J_{\lambda_i}\), and
\eqref{eq:support-coarse} gives
\begin{equation}
                  r_{\lambda_i}\geq c_*:=\ell_j/2>0
                       \quad\text{on }\supp_B F_i.          \label{eq:bounded-band-radius}
\end{equation}
The coordinates used by these radii escape every fixed initial
segment. Indeed, each coordinate occurs in at most two radii, so a
finite set of coordinates occurs in only finitely many levels.
After discarding finitely many terms, choose integers \(M_i\to\infty\)
so that every coordinate used by \(r_{\lambda_i}\) has index \(>M_i\).
Then \(r_{\lambda_i}(s)\leq\tau_{M_i}(s)\). Condition \eqref{eq:bounded-band-radius} is
the hypothesis of Lemma~\ref{lem:tail-separation}, which gives
detectors in one enlarged class with pairwise disjoint supports.
Lemma~\ref{lem:compatible}(1) gives a common Lipschitz bound for all
their finite signed sums.

\paragraph{Distinct labels with bands tending to infinity.}
Now take \(j_i\to\infty\). The fixed class must be \(\mathcal T_h\),
because a \(\mathcal V_h\) representation has \(j\leq h\).
By diagonal extraction from the compact intervals \([-1,1]\),
the finite tags \(\xi_i\in[-1,1]^{j_i}\) converge coordinatewise to
some \(v_*\in[-1,1]^\N\). This point need not be in \(c_0\).
Only its finite coordinate restrictions will be used.

Write \(a_i=2d_{j_i}+\ell_{j_i}/2\), which tends to zero.
For every fixed nonempty finite coordinate set \(I\), once
\(j_i\geq\max I\), \eqref{eq:support-coarse} gives
\begin{equation}
 \sup_{s\in\supp_B F_i}\|s|_I-v_*|_I\|_\infty
 \leq a_i+\|\xi_i|_I-v_*|_I\|_\infty\longrightarrow0.       \label{eq:projected-concentration}
\end{equation}
For each fixed \(i\), consider arbitrarily small positive dyadic
scales \(\theta\), below the last existing scale if there is one and
with \(\theta\) less than the remaining strict budget in
\eqref{eq:scale-rule}. Choose a grid center \(q_i(\theta)\) within
\(\theta/8\) of \(v_*|_{I_{\gamma_i}}\). The test
\[
                 F_{i,\theta}=G_{q_i(\theta),\theta}F_i
\]
belongs to \(\mathcal T_{h+1}\), as does \(F_i\), and
\[
       \|F_{i,\theta}-F_i\|_\infty
                    \leq\theta\Lip(F_i)\leq\theta4^hL_0.
\]
Thus \(F_{i,\theta}\to F_i\) pointwise as \(\theta\downarrow0\),
even though the centers move. For this fixed \(i\),
Lemma~\ref{lem:duality} and \eqref{eq:fixed-vector} allow a choice
\(\theta_i\) with
\[
 |\langle F_{i,\theta_i}-F_i,m_i\rangle|<\delta/2.
\]
This is a separate choice after fixing \(m_i\), not a uniform limit
over the sequence of vectors. Write \(\widetilde F_i=F_{i,\theta_i}\).
It detects \(m_i\) with bound \(\delta/2\).

The function \(\widetilde F_i\) is constant on the relatively open
cylinder
\begin{equation}
 C_i=\{s\in B:\|s|_{I_{\gamma_i}}-v_*|_{I_{\gamma_i}}\|_\infty
                                      <3\theta_i/8\},        \label{eq:plateau-cylinder}
\end{equation}
because the distance from its grid center is then less than
\(\theta_i/2\). Its enlarged representation still satisfies the
strict budget, so \eqref{eq:support-coarse} also applies to
\(\widetilde F_i\). Therefore \eqref{eq:projected-concentration} holds with \(\widetilde F_i\)
in place of \(F_i\). For each fixed earlier \(i\), all sufficiently
late supports lie in \(C_i\). A recursive subsequence makes this
true for every later selected index simultaneously. Now
Lemma~\ref{lem:compatible}(2), with \(u_i=\widetilde F_i\) and
\(b_i=1\), gives a common Lipschitz bound for finite signed sums.

\paragraph{One repeated label.}
It remains to treat a fixed label \(\gamma=(\lambda,j,\xi)\).
Here all cylindrical factors depend on the same finite set
\(I=I_\gamma\). We first expand the tests to isolate a varying term
that still detects the vectors. An escaping tail index then gives
disjoint supports. Otherwise, a first small remainder \(R\) gives
shrinking cylindrical supports; a final local modification produces
nested plateaux with summably small values.

Extract a subsequence on which the number \(m\leq h\)
and the types \(G,R\) of cylindrical operations are fixed, as are
the number \(l\leq h\) and the types of tail factors. For each
operation position its dyadic scale has a subsequence that is either
constant or tends to zero. In the constant case its center can also
be made constant, because the permitted grid in the fixed cube is
finite. At each tail position the scale \(c\) can be made constant,
and its integer index \(N\) is either constant or tends to infinity.
Make all these finitely many extractions together.

\emph{Expansion and detection.}
Call an ingredient fixed when all of its parameters are now constant.
The parameters in each concrete representation are frozen before
the following algebraic expansion. Replace every nonfixed
\(G_{q,\varepsilon}\) by
\(\mathrm{Id}+R_{q,\varepsilon}\), and every nonfixed factor
\(b_{c,N}\) by \(1+(b_{c,N}-1)\). The operations are linear in their
input functions, so these identities distribute through later
operations with no parameters recomputed. The expansion is only of
the cylindrical operators and of the tail product separately; it
never applies \(G\) to a tail-dependent full test.

There are at most \(K_0=2^{m+l}\leq2^{2h}\) resulting term positions.
Each cylindrical term retains a chronological subsequence of the
original operations, possibly replacing \(G\) by \(R\) at the same
scale. Its scales remain nonincreasing and its total scale does not
increase. Omitting a tail factor or replacing \(b\) by \(b-1\)
also preserves the permitted type and scale. Because \(j\) is fixed,
all resulting full terms belong to one class \(\mathcal V_H\) for
some fixed \(H\geq\max\{h,j\}\), even if the original class was
\(\mathcal T_h\). Write the exact expansion as
\begin{equation}
                             F_i=\sum_{a=1}^{K}F_i^a,
                 \qquad K\leq K_0.                          \label{eq:label-expansion}
\end{equation}

A term position containing only fixed ingredients is a fixed concrete
test \(F^a\). Its pairing with \(m_i\) tends to zero by
weak nullity; it is not asserted to vanish identically.
There are only finitely many such positions, so their total pairing
tends to zero. In view of \eqref{eq:selection-detection}, for all sufficiently large \(i\)
the sum of the absolute pairings of the varying terms in \eqref{eq:label-expansion} is
at least \(\delta/2\). One of the at most \(K_0\) positions has
pairing at least \(\delta/(2K_0)\); a subsequence fixes that position.
Relabel its tests as
\begin{equation}
              F_i=u_i b_i,\qquad
              |\langle F_i,m_i\rangle|\geq\delta_0>0.       \label{eq:varying-detection}
\end{equation}

If \(b_i\) contains a nonfixed factor \(b_{c,N_i}-1\) with
\(N_i\to\infty\), its support is contained in
\(\{\tau_{N_i}\geq c\}\) by \eqref{eq:tail-support}.
The band is fixed, so Lemma~\ref{lem:tail-separation} applies.
We may therefore assume no such factor remains. Every retained
nonfixed tail factor would have this form after expansion. The remaining
tail product is therefore fixed, and the variation in
\eqref{eq:varying-detection} must occur in the cylindrical factor.
Every retained nonfixed cylindrical operation is an \(R\): a
nonfixed \(G\) was replaced by an omitted identity or by \(R\).

\emph{A small remainder.}
Choose the first retained \(R\) whose scale
\(\varepsilon_i\to0\), and call its center \(q_i\).
All retained operations before it are fixed, so they produce one
fixed prefix function \(p\) from \(f_\gamma\). Set
\[
                 L_{\mathrm{pre}}=\Lip(p)\leq4^hL_0.
\]
This is the Lipschitz bound of the prefix, not the base bound \(L_0\);
earlier operations may have enlarged it. If \(L_{\mathrm{pre}}=0\),
this \(R\) and the resulting test vanish, contradicting \eqref{eq:varying-detection}.
Lemma~\ref{lem:operations} gives
\[
 \|R_{q_i,\varepsilon_i}p\|_\infty
       \leq L_{\mathrm{pre}}\varepsilon_i,\qquad
 K_I(R_{q_i,\varepsilon_i}p)
       \subset\overline B_\infty(q_i,\varepsilon_i).
\]
Every retained later scale is at most \(\varepsilon_i\), by the
nonincreasing order in the original representation. A later \(G\)
does not increase the supremum norm, a later \(R\) increases it
by at most a factor two, and either enlarges projected support by
at most its scale. Iterating the two parts of
Lemma~\ref{lem:operations} over at most \(h\) later operations yields
\begin{equation}
 \|u_i\|_\infty\leq2^hL_{\mathrm{pre}}\varepsilon_i,\qquad
 \{s|_I:s\in\supp_B u_i\}
    \subset K_I(u_i)
    \subset\overline B_\infty(q_i,(h+1)\varepsilon_i).       \label{eq:small-r-support}
\end{equation}
The first inclusion follows by continuity of coordinate restriction
and the definition of relative support. This is where the order of
scales is essential: a fixed total budget alone would not make the
sum of the later scales tend to zero.

Pass to a subsequence with \(q_i\to v\) in the compact cube
\([-1,1]^I\). We now make each cylindrical factor constant on a small
cylinder about \(v\), while preserving \eqref{eq:varying-detection}. For a fixed \(i\),
choose arbitrarily small positive dyadic \(\theta\) that are below
the last retained scale, below the remaining strict budget, and
at most \(\varepsilon_i\). Let \(q_i^+(\theta)\) be a permitted
grid point within \(\theta/8\) of \(v\), and set
\[
 \widetilde F_{i,\theta}
             =(G_{q_i^+(\theta),\theta}u_i)b_i.
\]
All these full tests, and \(F_i=u_i b_i\), belong to one fixed
enlarged class \(\mathcal V_{H'}\), with \(H'\) large enough for one
more cylindrical operation. The tail product has supremum norm at
most one, so
\[
 \|\widetilde F_{i,\theta}-F_i\|_\infty
       \leq\theta\Lip(u_i).
\]
For the fixed vector \(m_i\), pointwise convergence and
\eqref{eq:fixed-vector} give an allowed choice \(\theta_i\) such that
\begin{equation}
 |\langle\widetilde F_{i,\theta_i}-F_i,m_i\rangle|
                                      <\delta_0/2.           \label{eq:retained-detection}
\end{equation}
Again the scale is chosen separately for each \(m_i\); no uniform
dual convergence on the moving sequence is used.

Write \(\widetilde u_i=G_{q_i^+(\theta_i),\theta_i}u_i\) and
\(\widetilde F_i=\widetilde u_i b_i\). It detects with bound
\(\delta_0/2\). The support enlargement in
Lemma~\ref{lem:operations}, \eqref{eq:small-r-support}, and \(\theta_i\leq\varepsilon_i\)
give
\begin{equation}
 \{s|_I:s\in\supp_B\widetilde u_i\}
        \subset\overline B_\infty(q_i,(h+2)\varepsilon_i).  \label{eq:modified-small-support}
\end{equation}
This conclusion does not require the new grid center to be close to
\(q_i\): a new nonzero value can be created only near an old nonzero
value, which is exactly the support enlargement proved earlier.
As \(q_i\to v\) and \(\varepsilon_i\to0\), these cylindrical supports
concentrate at \(v\).

On the relatively open cylinder
\[
 C_i=\{s\in B:\|s|_I-v\|_\infty<3\theta_i/8\},
\]
the function \(\widetilde u_i\) is the constant
\(e_i=u_i(q_i^+(\theta_i))\). Equation \eqref{eq:small-r-support} gives
\begin{equation}
                   |e_i|\leq2^hL_{\mathrm{pre}}\varepsilon_i
                                \longrightarrow0.           \label{eq:small-plateau}
\end{equation}
\emph{Nested plateaux.}
Choose a recursive subsequence on which
\(\sum_i|e_i|<\infty\) and \(\supp_B\widetilde u_k\subset C_i\)
whenever \(k>i\). To see that both requirements can be imposed,
after selecting finitely many indices use \eqref{eq:modified-small-support} to discard a finite
initial segment so that all remaining supports lie in their
cylinders; then choose the next index with, for example,
\(|e_i|\leq2^{-i}\), using \eqref{eq:small-plateau}. Continue.

The remaining \(b_i\) are products of at most \(h\) tail factors of
supremum norm at most one and Lipschitz constant at most \(2^h\).
Thus \(\Lip(b_i)\leq h2^h\) uniformly. The full
\(\widetilde F_i\) have a common Lipschitz bound because they are
in \(\mathcal V_{H'}\), and
\[
                   \sum_i|e_i|\Lip(b_i)<\infty.
\]
Lemma~\ref{lem:compatible}(2), applied to the cylindrical supports
in \eqref{eq:modified-small-support}, gives the required common Lipschitz bound for all finite
signed sums. The nesting is of \(\supp_B\widetilde u_i\), not merely
the possibly smaller supports of \(\widetilde F_i\).

In each case the modified tests remain in one listed class, retain a
positive common lower bound on their pairings, and have a common
Lipschitz bound for finite signed sums. Choose an integer \(C_1\)
no smaller than this last bound. These are the required conclusions.
\end{proof}

\subsection{Excluding a linear copy of \texorpdfstring{\(c_0\)}{c0}}

\begin{proposition}\label{prop:no-c0}
The Banach space \(Z\) contains no linearly isomorphic copy of \(c_0\).
\end{proposition}

\begin{proof}
Suppose \(T:c_0\to Z\) is an isomorphic embedding and put \(m_i=Te_i\).
The lower bound for \(T\) and the identity
\(\|\sum_{i\in H}\epsilon_i e_i\|_\infty=1\) for nonempty finite
\(H\) give \eqref{eq:bounded-signed} for some \(r,M>0\).
Lemma~\ref{lem:detectors} supplies a weakly null subsequence and
detectors in one class. Proposition~\ref{prop:compatible-detectors}
then supplies compatible detectors with a common nonzero pairing.
For their class \(\mathcal L_1\), bound \(C_1\), and pairing
constant \(\delta_1\), set \(d_1=(\mathcal L_1,C_1,2)\).
Lemma~\ref{lem:detector-growth} gives
\[
                  M^2\geq\alpha_{d_1}|H|\delta_1^2
\]
for every finite set of selected indices \(H\), a contradiction.
Thus no such embedding \(T\) exists.
\end{proof}
